\documentclass[11pt,a4paper]{article}

\usepackage[margin=2.6cm]{geometry}
\usepackage[T1]{fontenc}
\usepackage{lmodern}
\usepackage{microtype}
\usepackage{amsmath,amssymb}
\usepackage{booktabs}
\usepackage{enumitem}

\setlength{\parindent}{0pt}
\setlength{\parskip}{6pt}
\setlist{nosep,leftmargin=1.6em}

\title{Initialization Suggestions for \texttt{misclassCRC}}
\author{Internal technical note}
\date{24 August 2026}

\begin{document}

\maketitle

\section{Purpose}

This note proposes an initialization scheme for the EM algorithm in
\texttt{misclassCRC}. It follows the structure of
\texttt{baffour\_misclass\_full()} in the accompanying simulation code while
generalizing the two-class random split to an arbitrary number of latent
classes. The initializer should consume the model matrices and aligned state
representation already constructed by \texttt{crc\_fit()}. Initialization ends
when the initial state weights and complete capture-cell counts have been
constructed.

\section{Notation}

For observed unit $i$, let $g$ denote a candidate configuration of the true
values of all misclassified variables and let $k\in\{1,\ldots,K\}$ denote a
latent class. Write
\begin{align*}
  M_i(g) &= \Pr(\widetilde G_i\mid G_i=g),\\
  q_{igk}^{(0)} &= \text{initial weight of state }(i,g,k),\\
  c(i,g,k) &= \text{complete capture cell associated with state }(i,g,k).
\end{align*}
When several variables are misclassified, $M_i(g)$ is the product of their
confusion-matrix contributions, consistently with the conditional-independence
assumption in the current package specification. When \texttt{true\_if}
establishes a true value, the corresponding contribution is replaced by an
indicator that the candidate value equals the observed value.

\section{Initial state weights}

For each admissible pair $(i,g)$, draw a mildly perturbed latent-class
allocation
\[
  (a_{ig1},\ldots,a_{igK})
  \sim \operatorname{Dirichlet}(\alpha,\ldots,\alpha).
\]
A moderately large value, such as $\alpha=20$, keeps the allocation close to
$1/K$ but breaks the exact symmetry between latent classes. The initialization
should be random whenever $K>1$; equal deterministic allocation can preserve
class symmetry and prevent the latent components from separating. The initial
state weights are
\[
  q_{igk}^{(0)}
  =
  \frac{M_i(g)a_{igk}}
       {\displaystyle\sum_{g'}\sum_{k'=1}^{K}
        M_i(g')a_{ig'k'}}.
\]
Thus, the weights are non-negative and sum to one within every observed unit.
Without misclassification, $g$ has only one admissible value and the formula
reduces to the perturbed latent-class allocation. With $K=1$, every admissible
state is weighted only through the misclassification mechanism.

The initializer should use R's current random-number generator state and should
not accept a seed argument. A user who requires reproducibility can call
\texttt{set.seed()} before \texttt{crc\_fit()} or before the internal
initializer is evaluated. This follows the usual R interface and avoids
changing the global random-number generator state inside the function.

\section{Choice of the Dirichlet concentration}

The symmetric Dirichlet concentration $\alpha$ controls the amount of random
variation around equal latent-class allocation. Each component has
\[
  \operatorname{E}(a_{igk})=\frac{1}{K},
  \qquad
  \operatorname{Var}(a_{igk})
  =\frac{K-1}{K^2(K\alpha+1)}.
\]
Values below one favor allocations near the corners of the simplex, whereas
large values concentrate the allocations around $1/K$. Initialization needs a
moderate concentration: allocations should be unequal enough to break class
symmetry but should not routinely create nearly empty classes.

The recommended package default is
\[
  \alpha=20.
\]
For $K=2$, the simulation code draws one class proportion uniformly between
$0.35$ and $0.65$. A symmetric beta distribution with the same variance has
$\alpha\approx16.17$, so the proposed default produces a similar, slightly
more concentrated random perturbation.

The value should be available as \texttt{control\$init\_alpha}, rather than as
a formal argument of \texttt{crc\_fit()}. The package should require it to be a
single, finite, strictly positive number. It is unused when $K=1$. The package
should not calculate $\alpha$ from a target variance: the fixed default keeps
the implementation transparent, while the control setting permits deliberate
sensitivity analysis. This quantity affects only the random starting values;
it is not a parameter of the fitted statistical model.

\section{Initial capture-cell counts}

For each complete capture cell $c$, aggregate the state weights as
\[
  n_c^{(0)}
  =
  \sum_i\sum_g\sum_{k=1}^{K}
  q_{igk}^{(0)}
  \mathbf{1}\{c(i,g,k)=c\}.
\]
The existing capture-cell index provides the required mapping without a new
join. For an all-zero capture cell, there is no observed contribution. A
simple initial rule, matching the simulation implementation, is
\[
  n_c^{(0)}=1
  \qquad\text{for every all-zero cell }c.
\]
These values are starting values rather than fixed pseudo-observations; the EM
updates subsequently replace them with model-based expected counts. This step
completes initialization.

\section{Recommended initialization sequence}

The initializer should perform exactly the following operations:
\begin{enumerate}
  \item Generate random Dirichlet allocations near equal class probabilities.
  \item Multiply them by the applicable misclassification contributions and
        normalize the resulting state weights within each observed unit.
  \item Aggregate state weights into initial expected capture-cell counts.
  \item Assign a starting expected count of one to every all-zero capture cell.
\end{enumerate}

The internal initialization object should retain the state weights, expected
cell counts, and initialization status. Its dimensions should be checked
against the existing state rows and capture-cell rows.

\section{Robustness extensions}

Because the initialization is random, repeated calls can produce different
starting state weights and capture-cell counts. Reproducible collections of
starting values can be generated by setting R's random-number seed once before
the repeated calls.

\section{Points to confirm during implementation}

The starting value of one for every all-zero cell and the randomized class
split reproduce the essential initialization choices in the simulation code.
Before exposing either choice through \texttt{control}, they should be checked
against the accompanying paper.

\end{document}
